\documentclass[11pt]{article}
\usepackage{fontspec}
\setmainfont{Latin Modern Roman}
\usepackage{amsmath,amssymb,mathtools,unicode-math}
\setmathfont{Latin Modern Math}
\usepackage{xurl,hyperref}
\hypersetup{colorlinks=true,urlcolor=blue,pdftitle={Problem statement}}
\newcommand{\sref}[2]{\hyperlink{source:#1:#2}{[S#2]}}
\newcommand{\cataloguescope}[1]{\paragraph{Mathematical question.}#1}
\begin{document}
\section{Algebraic Stories — Lefschetz properties of graded algebras}\label{statement}

\subsection{Definitions and mathematical statement}

\paragraph{Shared notation.}$\mathbb N=\{1,2,\ldots\}$, and $\mathbb Z,\mathbb Q,\mathbb R,\mathbb C$ denote the integers, rationals, reals and complex numbers. Any other conventions are specified locally.


For a graded algebra $A=\bigoplus_{i\geq0}A_i$, multiplication by a generic linear form $\ell\in A_1$ has maximal rank if its rank is the minimum of the dimensions of its domain and codomain. The weak Lefschetz property (WLP) requires maximal rank for $\ell:A_i\to A_{i+1}$ for all $i\geq0$. The strong Lefschetz property (SLP) requires maximal rank for $\ell^e:A_i\to A_{i+e}$ for every $i\geq0$ and $e\geq1$.For $n,d,k\geq1$, put
\[T_{n,d,k}=\mathbb C[x_1,\ldots,x_n]/(x_1^d,\ldots,x_n^d)^k.
\]
The exponent $k$ denotes the ordinary ideal power.
\paragraph{Statement recorded in the supplied source.}
Characterize all triples $(n,d,k)$ for which $T_{n,d,k}$ has the WLP, and all triples for which it has the SLP. \sref{statement}{2}
\subsection{Short English statement}
Determine exactly when the stated powers of monomial complete-intersection ideals give weak or strong Lefschetz algebras.
\subsection{Sources}
\begin{description}
\item[\hypertarget{source:statement:1}{S1}] ULAM.AI and the UnsolvedMath Contributors, UnsolvedMath: A Curated Collection of Open Mathematics Problems, catalogue entry. CC BY 4.0. \url{https://huggingface.co/datasets/ulamai/UnsolvedMath}.
\item[\hypertarget{source:statement:2}{S2}] R. Fröberg, S. Lundqvist, A. Oneto and B. Shapiro, \emph{Algebraic Stories from One and from the Other Pockets}, arXiv:1801.01692 (2018), §2.3, Problem H. \url{https://arxiv.org/abs/1801.01692}.
\end{description}
\label{end:statement}
\end{document}
